\recordintro{概要}
\notelabel{block:talk01-overview}

この節には，講演または集中講義の最初の話題ブロックを置く。
完成 PDF では，作業資料名やページ番号を本文に出さず，読者がこの PDF だけで読める形で本文を書く。

本稿を AI と共同で作成した場合は，ノート作成者が講義中に作成した手書きノートを主資料としていること，
AI を補助的に使用したこと，使用した AI の名称，モデル名，推論設定を実際の使用状況に合わせて明記すること，
内容の最終確認および誤りの責任がノート作成者にあることを概要に明記する。

参考文献との対応が確定している主張には，必要に応じて本文中で引用を付ける。
未確認の引用候補や判読不能箇所は，完成 PDF へ作業メモとして出さず，作業台帳に分離する。

\recorddate{YYYY年M月D日}
\section{最初の講義日}

この節には，その日の講義内容を手書きノートの順序に沿って置く。
日付は記録として本文に表示するが，目次には入れない。
